\begin{abstract}
    Questo report analizza la discrepanza strutturale nella misurazione dell'emigrazione italiana, focalizzandosi sul sottoconteggio sistematico dei cittadini espatriati, rilevata dall'ISTAT.
    Attraverso un approccio di \textit{statistiche speculari}, abbiamo confrontato i dati ufficiali ISTAT (flussi in uscita) con i record di immigrazione Eurostat (flussi in entrata) per otto paesi europei nel periodo 2017-2021.
    L'analisi integra due test non parametrici, ovvero, quello di \textit{Wilcoxon} per valutare l'entità del bias e il $\chi^2$ per mappare la deviazione della distribuzione temporale.
    I risultati dimostrano con forte evidenza statistica una sottostima da parte dell'ISTAT, distorcendo così in modo critico la rilevazione degli italiani non più residenti nel territorio.
\end{abstract}